\documentclass{article}
\usepackage[UTF8]{ctex}
\usepackage{amsmath,amssymb,amsfonts}
\usepackage{enumitem}
\usepackage{geometry}
\geometry{a4paper, margin=1in}
% === 错题本排版优化（不参与题目标识） ===
\setlength{\parindent}{0pt}
\setlength{\parskip}{0.35em plus 0.10em minus 0.08em}
\setlength{\emergencystretch}{3em}
\setlist{itemsep=0.25em, topsep=0.40em, parsep=0pt}
\allowdisplaybreaks[2]
% === 错题本排版优化结束 ===
\begin{document}
\section*{华中师范大学2026年数学分析真题}
\section*{140．华中师范大学2026年数学分析试卷（回忆版）}
计算题．

\begin{itemize}
  \item[（1）] 求极限 \(\lim _{x \rightarrow 0} x^{-4}\left(\cos x-e^{-\frac{x^{2}}{2}}\right)\) ．
  \item[（2）] 求由 \(z=x^{2}+y^{2}, z=x+y\) 所围成的立体图形的体积。
  \item[（3）] 计算曲面积分
\[
-\iint_{S} x z^{2} \mathrm{~d} y \mathrm{~d} z+\left(x^{2} y-z^{3}\right) \mathrm{d} z \mathrm{~d} x+\left(2 x y+y^{2} z\right) \mathrm{d} x \mathrm{~d} y .
\]
其中 \(S\) 是由 \(z=\sqrt{a^{2}-x^{2}-y^{2}}(a>0)\) 与 \(z=0\) 所围立体的表面，取外侧为正向．
\end{itemize}

\subsection{题目1}
二．解答如下问题：

\begin{itemize}
  \item[（1）] 设 \(L\) 为常数，\(\left\{y_{n}\right\}\) 是严格递增的正无穷大数列，且与数列 \(\left\{x_{n}\right\}\) 一起满足 \(\lim _{n \rightarrow \infty} \frac{x_{n+1}-x_{n}}{y_{n+1}-y_{n}}=L\) 。证明： \(\lim _{n \rightarrow \infty} \frac{x_{n}}{y_{n}}=L\) ．
\end{itemize}

\subsection{题目1}
原题边界“（2）”：设 \(a_{1}=\sin a>0, a_{n+1}=\sin a_{n}(n=1,2, \cdots)\) ，证明： \(\lim _{n \rightarrow \infty} \frac{1}{n a_{n}^{2}}=\frac{1}{3}\) ．

\subsection{题目2}
三．设 \(f(x)\) 在 \((a, b)\) 上可导， \(\lim _{x \rightarrow a^{+}} f^{\prime}(x)\) 与 \(\lim _{x \rightarrow b^{-}} f^{\prime}(x)\) 均存在．证明：

\begin{itemize}
  \item[（1）] \(f(a+0)\) 和 \(f(b-0)\) 都存在．
  \item[（2）] \(f(x)\) 在 \((a, b)\) 上有界，且一致连续．
\end{itemize}

\subsection{题目3}
四．设 \(f(x)\) 在 \(I=(0,+\infty)\) 上可导，且对任意的 \(x_{1}, x_{2} \in I\) 以及 \(\lambda_{1}, \lambda_{2} \in(0,1)\) ，有

\[
f\left(\lambda_{1} x_{1}+\lambda_{2} x_{2}\right) \leq \lambda_{1} f\left(x_{1}\right)+\lambda_{2} f\left(x_{2}\right) .
\]

其中 \(\lambda_{1}+\lambda_{2}=1\) 。假设 \(x \in(0,+\infty)\) 时，有 \(f(x)=x^{2}+x^{2} \varepsilon(x)\) ，其中 \(\lim _{x \rightarrow+\infty} \varepsilon(x)=0\) 。证明：

\begin{itemize}
  \item[（1）] 对任意的 \(x \in(0,+\infty)\) ，有 \(\frac{f(x)-f(x-h)}{h} \leq f^{\prime}(x) \leq \frac{f(x+h)-f(x)}{h}, 0<h<\frac{x}{2}\) ．
  \item[（2）] 对任意的 \(\eta>0\) ，当 \(x\) 充分大时，有 \(2 x-h-\frac{\eta}{h} x^{2} \leq f^{\prime}(x) \leq 2 x+h+\frac{\eta}{h} x^{2}, 0<h<\frac{x}{2}\) ．
\end{itemize}

\subsection{题目4}
五．设 \(f(x)\) 为有界闭区间 \([a, b]\) 上的函数，且对任意的 \(x_{0} \in[a, b]\) ，有 \(\lim _{x \rightarrow x_{0}} f(x)=0\) ．证明：

\begin{itemize}
  \item[（1）] \(f(x)\) 在 \([a, b]\) 上有界．
  \item[（2）] 对任意的 \(\varepsilon>0\) ，在 \([a, b]\) 上至多存在有限个点 \(x\) ，使得 \(|f(x)| \geq \varepsilon\) ．
  \item[（3）] \(f(x)\) 在 \([a, b]\) 上可积，且 \(\int_{a}^{b} f(x) \mathrm{d} x=0\) ．
\end{itemize}

\subsection{题目5}
六．判断 \(\sum_{n=3}^{\infty} \frac{1}{n^{p}(\ln n)^{q}(\ln \ln n)^{r}}\) 的敛散性。

\subsection{题目6}
七．设函数列 \(\left\{f_{n}(x)\right\}\) 与 \(\left\{g_{n}(x)\right\}\) 在区间 \(I\) 上分别一致收敛于 \(f(x)\) 与 \(g(x)\) ，且存在正数列 \(\left\{M_{n}\right\}\) ，使得当 \(x \in I, n=1,2, \cdots\) 时，有 \(\left|f_{n}(x)\right| \leq M_{n},\left|g_{n}(x)\right| \leq M_{n}\) 。证明：\(\left\{f_{n}(x) g_{n}(x)\right\}\) 在 \(I\) 上一致收敛。

\subsection{题目7}
八．已知 \(f(x)\) 是以 \(2 \pi\) 为周期的可积函数，通过 \(f(x)\) 的傅里叶系数表示 \(f(x+h)\) 的傅里叶系数．

\subsection{题目8}
九．设 \(f(x, y)=|x-y| g(x, y)\) ，其中 \(g(x, y)\) 在 \((0,0)\) 的某邻域内连续．问：

\begin{itemize}
  \item[（1）] 在什么条件下，\(f_{x}(0,0), f_{y}(0,0)\) 存在？
  \item[（2）] 在什么条件下，\(f(x, y)\) 在 \((0,0)\) 处可微？
\end{itemize}

\subsection{题目9}
十．设函数 \(f(x)\) 在 \((0,+\infty)\) 上有界且连续，若对任意的 \(c \in \mathbb{R}\) ，函数 \(f(x)-c\) 在 \((0,+\infty)\) 上至多有有限个零点，证明： \(\lim _{x \rightarrow+\infty} f(x)\) 存在．
\end{document}
